\pageheading{ORBIT SPECTRA\; /\; 1.0 THEOREM}{The continuation gap is closed}
For every integer $N\ge2$, the orbit zeta has a meromorphic continuation to
the whole complex plane with exactly one pole:
\[
\boxed{s=1/N,\qquad
\operatorname{Res}_{s=1/N}\mathcal Z_N(s)=\frac{\sin(\pi/N)}{N\pi}.}
\]
Its exact absolute-convergence abscissa is $1/N$. At nonpositive integers,
\[
\boxed{\mathcal Z_N(-m)=0\quad(m\ge0),\qquad
\mathcal Z_N'(0)=-\log N.}
\]
The zeta-regularized product of the orbit coordinates is consequently
$\exp[-\mathcal Z_N'(0)]=N$.
\fig{15_spectral_poles}{Pole and zero locations follow from the proved
continuation. The second panel shows the shifted extension, where a pole
ladder appears. These are exact structural locations rather than fitted roots.}
The continuation uses classical Mellin regularization. The contribution at
this checkpoint is a complete local application with multiplicities,
principal branches, and all convergence hypotheses checked.

\pageheading{ORBIT SPECTRA\; /\; LARGE RADIUS}{A finite exponential sum controls the tail}
Put $E_N(t)=d_N(-t)=\sum_{n\ge0}t^n/(Nn)!$ and
$\omega=e^{2\pi i/N}$. For $x\ge0$,
\[
E_N(x^N)=\frac1N\sum_{j=0}^{N-1}e^{\omega^jx}.
\]
Root filtering proves this identity term by term: only powers divisible by
$N$ remain. The positive real exponential is dominant, so
\[
\log E_N(t)=t^{1/N}-\log N+O(e^{-\kappa_N t^{1/N}}),
\quad \kappa_N=1-\cos(2\pi/N)>0,
\]
with corresponding differentiated estimates on the positive ray.
\fig{14_spectral_tail}{Computed logarithmic derivatives of the finite
exponential sum approach one in the root-radius variable. The decaying
remainder supplies the Mellin continuation's large-radius input.}
\claim{Why complex zero arguments do not break the integral}
In a sufficiently narrow sector around negative $w$, choose the root
$x=(-w)^{1/N}$ near the positive real ray. The first exponential dominates
the other $N-1$ terms by a uniform real-part gap at large radius, so this
sector is eventually zero free. The finitely many remaining zeros avoid
the negative real ray because $d_N(-t)>0$ for $t\ge0$.
Thus all zero arguments have a uniform gap from the branch cut.
In particular $|t+w_j|\ge c_N(t+|w_j|)$ for $t\ge0$.

\pageheading{ORBIT SPECTRA\; /\; MELLIN PROOF}{The two ends have different jobs}
For $1/N<\Re s<1$, the genus-zero log derivative and the angular gap give
absolute interchange in
\[
\mathcal Z_N(s)=\frac{\sin\pi s}{\pi}
\int_0^\infty t^{-s}\frac{d}{dt}\log E_N(t)\dd t.
\]
Each term uses the principal-branch beta integral
$\int_0^\infty t^{-s}(t+w)^{-1}\dd t=\pi w^{-s}/\sin\pi s$.
The denominator estimate bounds the absolute integral by a constant times
$|w|^{-\Re s}$, whose sum converges in this strip.

Near zero write $f(t)=(\log E_N)'(t)=\sum_{k\ge1}\alpha_kt^{k-1}$,
where $\alpha_k=(-1)^{k+1}\mathcal Z_N(k)$.
For any integer $K\ge1$, subtract $K$ terms on $[0,1]$:
\[
\begin{aligned}
I_K(s)={}&\int_0^1t^{-s}\left[f(t)-\sum_{k=1}^K\alpha_kt^{k-1}\right]\dd t
+\sum_{k=1}^K\frac{\alpha_k}{k-s}\\
&+\int_1^\infty t^{-s}\left[f(t)-\frac1N t^{1/N-1}\right]\dd t
+\frac{1}{N(s-1/N)}.
\end{aligned}
\]
The first integral is holomorphic for $\Re s<K+1$; the second is entire
because its remainder decays exponentially in $t^{1/N}$.
The factor $\sin\pi s$ cancels the small-radius poles at positive integers.
Letting $K$ increase covers the plane and leaves only the pole at $1/N$.
At nonpositive integers $I_K$ is regular, giving the stated zero values.

At zero the finite parts combine to
\[
I_K(0)=\lim_{T\to\infty}[\log E_N(T)-T^{1/N}]=-\log N.
\]
Since $(\sin\pi s/\pi)'_{s=0}=1$, this is $\mathcal Z_N'(0)$.
The pole also proves the exact abscissa: absolute convergence at any smaller
real exponent would give a holomorphic series across $1/N$, a contradiction.

\pageheading{ORBIT SPECTRA\; /\; SHIFTS}{A Hurwitz-type deformation and its determinant}
For real $a\ge0$, define initially
\[
\mathcal Z_N(s;a)=\sum_j(w_j+a)^{-s},\qquad\Re s>1/N.
\]
The rightward shift preserves avoidance of the negative branch cut.
Normalize the product with $E_N(a)$ and use
$\log E_N(a+t)$ in the Mellin argument.
For $a>0$ there are genuine simple poles at $s=1/N-k$, $k\ge0$:
\[
\boxed{\operatorname{Res}_{s=1/N-k}\mathcal Z_N(s;a)
=\frac{\sin(\pi/N)}{N\pi}\frac{(1-1/N)_k}{k!}a^k.}
\]
Indeed the large-radius derivative expands as
\[
\frac1N(a+t)^{1/N-1}
=\frac1N\sum_{k\ge0}\frac{(-1)^k(1-1/N)_k}{k!}a^kt^{1/N-1-k}
\]
as an asymptotic series. The alternating sign is canceled by
$\sin(\pi(1/N-k))=(-1)^k\sin(\pi/N)$.
At $a=0$ all secondary residues vanish.

The negative-integer values remain zero, and
\[
\mathcal Z_N'(0;a)=-\log[N E_N(a)],\qquad
\boxed{\prod_j^{\rm reg}(w_j+a)=N E_{N,1}(a).}
\]
This uses the large-radius finite part $-\log N$ and the normalization
term $-\log E_N(a)$. It is a zeta-regularized product, not an ordinary
convergent product of the coordinates.

For $N=2$, the result is $2\cosh\sqrt a$, a classical half-integer
spectral determinant. The four shifted numerical benchmarks reproduce this
anchor. Functional equations and derivative-weighted zero sums remain
separate questions; the continuation proof does not settle them.
